% lmcs, 2018.5.2, 2019.9.9, 2020.2.10


\documentclass{lmcs}


%\usepackage{fouriernc}
%\setmathfont[FakeBold=0.5]{Latin Modern Math}

\usepackage{wasysym}
\usepackage{hyperref}
\usepackage{graphicx}


\usepackage{amssymb,amsmath,graphicx,color}
\usepackage{mymath,proof,latexsym}
\usepackage{xcolor}
\usepackage[pdftex,outline]{contour}


%%%%%%%%%%%%%%%%%%%%%%%%%%%%%%%%%%%%
% WE ADD THE COMMANDS 
% theorem lemma proposition corollary definition  
%%%%%%%%%%%%%%%%%%%%%%%%%%%%%%%%%%%%

\newtheorem{theorem}{Theorem}[section]
\newtheorem{lemma}[theorem]{Lemma}
\newtheorem{proposition}[theorem]{Proposition}
\newtheorem{corollary}[theorem]{Corollary}
\newtheorem{definition}[theorem]{Definition}

%\newenvironment{proof}[1][Proof]{\begin{trivlist}
%\item[\hskip \labelsep {\bfseries #1}]}{\end{trivlist}}

\newenvironment{example}[1][Example]{\begin{trivlist}
\item[\hskip \labelsep {\bfseries #1}]}{\end{trivlist}}
\newenvironment{remark}[1][Remark]{\begin{trivlist}
\item[\hskip \labelsep {\bfseries #1}]}{\end{trivlist}}

%\bibliographystyle{plainurl}% the mandatory bibstyle for CSL2024

\begin{document}
\sloppy 
\hbadness=10000
\vbadness=10000

\def\cal{\mathcal} % for LMCS
\def\calH{{\cal H}}

% journal, pa2

\def\PAtwoID{{\rm PA}_2{\rm ID}}
\def\CPAtwoID{{\rm CPA}_2{\rm ID}}
\def\AxiomPA{{\rm AxiomPA}}

% cmcs

\def\Co{{\rm co}}
\def\Axiom{{\rm Axiom}}
\def\Wk{{\rm Wk}}
\def\Cut{{\rm Cut}}
\def\BS{{\rm BS}}
\def\Bit{{\rm Bit}}
\def\LK{{\rm LK}}
\def\Head{{\rm head}}
\def\Tail{{\rm tail}}

% equivalence cyclic proofs

%\def\Uor{\displaystyle \mathop {\widetilde\bigvee}}
\def\Uor{\biguplus}
\def\CLKIDomega{{{\bf CLKID}^\omega}}
\def\CLKIDPA{{{\bf CLKID}^\omega+{\bf PA}}}
\def\LKIDPA{{{\bf LKID}+{\bf PA}}}
\def\PA{{\bf PA}}

% equiv note

\def\Choose{{\rm Choose}}
\def\Asym{{\rm Asym}}
\def\Peano{{\bf Peano}}
\def\Code#1{{\lceil{#1}\rceil}}
\def\Ineq{{\rm Ineq}}
\def\MonoPath{{\rm MonoPath}}
\def\MonoSeq{{\rm MonoSeq}}
\def\BoundedPath{{\rm BoundedPath}}
\def\InfPath{{\rm InfPath}}
\def\InfMonoPath{{\rm InfMonoPath}}
\def\InfMonoSeq{{\rm InfMonoSeq}}
\def\HomSeq{{\rm HomSeq}}
\def\InfHomSeq{{\rm InfHomSeq}}

\def\ErdosTree{{\rm ErdosTree}}
\def\KB{{\rm KB}}
\def\MonoList{{\rm MonoList}}
\def\Seq{{\rm Seq}}
\def\Univ{{\rm univ}}
\def\Infinite{{\rm Infinite}}
\def\S{{\bf S}}
\def\N{{\bf N}}
\def\Ind{{\rm Ind}}
\def\LKID{{\bf LKID}}
\def\PA{{\bf PA}}
\def\LKIDExt{{\bf LKIDExt}}
\def\CLKID{{\bf CLKID}}

% LICS

\def\Over#1#2{\deduce{#2}{#1}}
\def\List{{\rm List}}
\def\ListX{{\rm ListX}}
\def\ListO{{\rm ListO}}
\def\ListE{{\rm ListE}}
\def\Nl{{\rm nl}}

% CSLID

\def\SLLID{{\rm SLID}}
\def\SLID{{\rm SLID}}
\def\CSLID{{\rm CSLID}}
\def\CSLIDomega{{{\rm CSLID}\omega}}
\def\Eclass{{\rm Eclass}}
\def\Eq{{\rm Eq}}
\def\Deq{{\rm Deq}}
\def\Satom{{\rm Satom}}
\def\Cells{{\rm Cells}}
\def\Roots{{\rm Roots}}
\def\Elim{{\rm Elim}}
% \def\Case#1{{\rm Case-$#1$}}
\def\Case{{\rm Case}}
\def\Unfold{{\rm Unfold}}
\def\Pred{{\rm Pred}}
\def\Start{{\rm Start}}
\def\Unsat{{\rm Unsat}}
\def\Split{{\rm Split}}
\def\Underscore{\underline{\phantom{x}}}
\def\Rootcell{{\rm rootcell}}
\def\Rootshape{{\rm rootshape}}
\def\Jointleaf{{\rm jointleaf}}
\def\Jointleafimplicit{{\rm jointleafimplicit}}
\def\Jointnode{{\rm jointnode}}
\def\Directjoint{{\rm directjoint}}

% cyclicnote2

\def\Range{{\rm Range}}
\def\Xstart{X_{{\rm start}}}
\def\kstart{k_{{\rm start}}}
\def\Leaves{{\rm Leaves}}
\def\PureDist{{\rm PureDist}}
\def\Pure{{\rm Pure}}
\def\Identity{{\rm Identity}}
\def\Subst{{\rm Subst}}

% weak establish

\def\Connected{{\rm Connected}}
\def\Eststablished{{\rm Eststablished}}
\def\Valued{{\rm Valued}}

% monadic new

\def\SLRDbtw{{\rm SLRD}_{btw}}
\def\BTW{{{\rm BTW}}}
\def\Loc{{{\rm Loc}}}
\def\Ne{{\ne}}
\def\Nil{{{\rm nil}}}
\def\eqDef{=_{{\rm def}}}
\def\Inf#1{{\infty_{#1}}}
\def\Stores{{\rm Stores}}
\def\SVars{{{\rm SVars}}}
%\def\Val{{{\rm Val}}}
\def\MSO{{{\rm MSO}}}
\def\Sep{{{\rm Sep}}}
\def\Septwo{{{\rm SLMI}}}
\def\SLMI{{{\rm SLMI}}}
\def\Sepinf{{\rm Sep}\infty}
\def\THeaps{{\rm THeaps}}
\def\Root{{\rm Root}}
\def\TG{{\rm TGraph}}

% monadic

\def\nil{{{\rm nil}}}
\def\eqDef{=_{{\rm def}}}
\def\Inf#1{{\infty_{#1}}}
\def\Stores{{\rm Stores}}
\def\SVars{{{\rm SVars}}}
\def\Val{{{\mathcal V}}}
\def\MSO{{{\rm MSO}}}
\def\Sep{{{\rm sep}}}
\def\THeaps{{\rm THeaps}}
\def\Root{{\rm Root}}
\def\TG{{\rm TGraph}}

\def\Tilde{\widetilde}
\def\Bar{\overline}
\def\Lequiv{\Longleftrightarrow}
\def\Lto{\Longrightarrow}
\def\Lfrom{\Longleftarrow}
\def\Norm{{\rm Norm}}
\def\Noshare{{\rm Noshare}}
\def\Roots{{\rm Roots}}
\def\Forest{{\rm Forest}}
\def\Var{{\rm Var}}
\def\Range{{\rm Range}}
\def\Cell{{\rm Cell}}
\def\Tree{{\rm Tree}}
\def\Switch{{\rm Switch}}
\def\All{{\rm All}}
\def\To{\leadsto}
\def\tree{{\rm tree}}
\def\Paths{{\rm Paths}}
\def\Finite{{\rm Finite}}
\def\Const{{\rm Const}}
\def\Leaf{{\rm Leaf}}
\def\LeastElem{{\rm LeastElem}}
\def\LeastIndex{{\rm LeastIndex}}
\def\WSnS{{{\rm WSnS}}}
\def\Expand{{{\rm Expand}}}

% previous

\long\def\J#1{} % skip
\def\T#1{\hbox{\color{green}{$\clubsuit #1$}}}
\def\W#1{\hbox{\color{Orange}{$\spadesuit #1$}}}

\def\Node{{{\rm Node}}}
\def\LL{{{\rm LL}}}
\def\DSN{{{\rm DSN}}}
\def\DCL{{{\rm DCL}}}
\def\LS{{{\rm LS}}}
\def\Ls{{{\rm ls}}}

\def\FPV{{{\rm FPV}}}
\def\Lfp{{{\rm lfp}}}
\def\IsHeap{{{\rm IsHeap}}}

\def\Equiv{\quad \equiv\quad }
\def\Null{{{\rm null}}}
\def\Emp{{{\rm emp}}}
\def\If{{{\rm if\ }}}
\def\Then{{{\rm \ then\ }}}
\def\Else{{{\rm \ else\ }}}
\def\While{{{\rm while\ }}}
\def\Do{{{\rm \ do\ }}}
\def\Cons{{{\rm cons}}}
\def\Dispose{{{\rm dispose}}}
\def\Vars{{{\rm Vars}}}
\def\Locs{{{\rm Locs}}}
\def\States{{{\rm States}}}
\def\Heaps{{{\rm Heaps}}}
\def\FV{{{\rm FV}}}
\def\True{{{\rm true}}}
\def\False{{{\rm false}}}
\def\Dom{{{\rm Dom}}}
\def\Abort{{{\rm abort}}}
\def\New{{{\rm New}}}
\def\W{{{\rm W}}}
\def\Pair{{{\rm Pair}}}
\def\Lh{{{\rm Lh}}}
\def\lh{{{\rm lh}}}
\def\Elem{{{\rm Elem}}}
\def\EEval{{{\rm EEval}}}
% \def\BEval{{{\rm BEval}}} % not used 
\def\PEval{{{\rm PEval}}}
\def\HEval{{{\rm HEval}}}
\def\EVal{{{\rm Eval}}}
\def\Domain{{{\rm Domain}}}
\def\Exec{{{\rm Exec}}}
\def\Store{{{\rm Store}}}
\def\Heap{{{\rm Heap}}}
\def\Storecode{{{\rm Storecode}}}
\def\Heapcode{{{\rm Heapcode}}}
\def\Lesslh{{{\rm Lesslh}}}
\def\Addseq{{{\rm Addseq}}}
\def\Separate{{{\rm Separate}}}
\def\Result{{{\rm Result}}}
\def\Lookup{{{\rm Lookup}}}
\def\ChangeStore{{{\rm ChangeStore}}}
\def\ChangeHeap{{{\rm ChangeHeap}}}
\def\Wand{\mathbin{\hbox{\hbox{---}$*$}}}
\def\Eval#1{[\kern -1pt[{#1}]\kern -1pt]}
\def\Vec{\overrightarrow}
%\def\Vec#1{{\bf #1}}

\def\Tilde{\widetilde}
\def\Break{\hfil\break\hbox{}}

\newcommand{\freshcopy}           [1]                      {  {\tilde{#1}} }
%\newcommand{\freshcopy}           [1]                      {  \displaystyle{\mathop{#1}^{\sim}} }
\newcommand{\bfColor} [2]     {{\bf \color{#1}{#2}}}
\newcommand{\arity}               {{\tt ar}}          
\newcommand{\trace}               {{\tt tr}} 
\newcommand{\stat}                 {{\tt stat}} 
\newcommand{\Stat}                {{\tt S}}     
\newcommand{\prog}               {{\tt p}} 
\newcommand{\Prog}               {{\tt P}}   
\newcommand{\fun}                 {{\tt fun}}          
\newcommand{\quotationMarks} [1]  {{\textquotedblleft #1\textquotedblright}}
\newcommand{\universe} [1]  {{| #1 |}}
\newcommand{\Rule}               {{\tt Rule}} 
\newcommand{\Path}               {{\tt Path}} 
\newcommand{\pathBetween} {{\tt path}} 
\newcommand{\Trace}               {{\tt Trace}} 
\newcommand{\Cyclic}              {{\tt Cyclic}}
\newcommand{\Cycle}              {{\tt Cycle}}
\newcommand{\dom}              {{\tt dom}}
\newcommand{\codom}              {{\tt codom}}
\newcommand{\Periodic}              {{\tt Periodic}}
\newcommand{\setR}              {{\mathcal R}}
\newcommand{\setL}              {{\mathcal L}}
\newcommand{\setG}              {{\mathcal G}}
\newcommand{\Pfin}      {{ {\mathcal R}_{\mbox{\tiny fin}} }}
\newcommand{\Card}              {{\tt Card}}
\newcommand{\restr}               {{\lceil}}
\newcommand{\comp}              {{\circ}}
\newcommand{\val}                   {{\tt val}}
\newcommand{\setP}                   {{\mathcal P}}
\newcommand{\Rprog}               {{\star}}
\newcommand{\RProg}               {{\bigstar}}
\newcommand{\length}               {{\mbox{\tt len}}}
\newcommand{\lex}                    {{\mbox{\tiny lex}}}
\newcommand{\setH}                  {{\mathcal H}}
\newcommand{\setS}                  {{\mathcal S}}
\newcommand{\setT}                  {{\mathcal T}}
\newcommand{\id}                      {{\tt id}}
\newcommand{\BrowerKleene}  {{\mbox{\tiny BK}}}
\newcommand{\conc}                 {{*}}





\title{Size Change Termination Condition 
for Injective Relations}
\date{}

\author[S. Berardi]{Stefano Berardi}
\address{Universit\`{a} di Torino,
Torino, Italy}
\email{stefano@unito.it}

%% mandatory lists of keywords 
\keywords{
proof theory
inductive definitions
cyclic proofs
infinite Ramsey theorem
}




\maketitle

\section{Introduction}
We want to study Neil Jones Size Change Termination condition \cite{2001-SCT}
(from now on, SCT) for program termination in a particular case: 
the case of Injective bi-colored Relations.

Our motivation is that this particular case is enough to treat a relevant
application of SCT, namely checking  the correctness of a cyclic proof.

We suspect that the algorithm checking SCT is 
PSPACE-complete also in this case. 
However we will argue that checking SCT is simpler  
and faster for injective relations.

Another motivation we have in mind is an automatic translation from 
cyclic proofs into more ordinary inductive proofs. 
We prove termination for SCT in the case of injective relations only using 
first order induction. With respect to a previous attempt,
we claim that we can translate cyclic proofs avoiding in the resulting proof
a formal proof of Ramsey Theorem and using Intuitionistic Logic only.


\section{Size Change Termination in the Case of Injective Relations}
%14:03 30/10/2024 
Along this paper we suppose we have: 
\begin{enumerate}
\item
a \emph{finite set} $X=\{x_1, \ldots, x_n\}$;
\item
a finite set $\setR =\{(\Stat_i, \Prog_i)| i \in [1,r]\}$, with
each $(\Stat_i, \Prog_i)$ a \emph{pair of finite relations on $X$}.

\item
a relation $\setG$ on 
$\setR \times (G \times G)$
\end{enumerate}

We call each $(\Stat_i, \Prog_i)$ just a \emph{pair} for short.

We define composition of pairs in $\setR$ by:
$$
(\Stat,\Prog)(\Stat',\Prog') 
=
(\Stat \Stat', \Stat \Prog' \cup \Prog \Stat' \cup \Prog \Prog')
$$
We motivate this definition later. For the moment, we stress that
we use $\setG$ to restrict the compositions which are allowed.

\begin{definition}[Composable pairs  in $\setG$]
Assume $\setR, \setG$ are as above.
\begin{enumerate}
\item
We say that $(\Stat,\Prog), (\Stat',\Prog') \in \setR$
are \emph{composable in $\setG$}, just \emph{composable}
for short, if and only if for some $a,b,c \in G$ we have 
$(\Stat,\Prog) \ \setG \ (a,b)$ and $(\Stat',\Prog') \ \setG \ (b,c)$
\item
A sequence of pairs of $\setR$ is
\emph{composable in $\setG$}, just \emph{composable} for short,
if and only if each pair is composable with the next pair.
\end{enumerate}
\end{definition}

 We consider the elements of $X$
as variables to which we associate a value in each instant of time.


We will use each relation $a \Stat_i b$ to express the fact that $a$
is associated in an instant $t$
to a value \emph{greater than or equal} to the value associated to $b$ in the next
instant $t+1$.
We will use each relation $a \Prog_i b$ to express the fact that $a$
is associated in an instant $t$
to a value \emph{greater than} the value associated to $b$ in the next
instant $t+1$.
% disjointness of \Stat_i and \Prog_i is NOT REQUIRED
We call \emph{trace} any sequence on $X$ 
and \emph{path} any sequence of pairs of $\setR$
which is \underline{composable} in $\setG$.

We take the words \quotationMarks{trace} and \quotationMarks{path} 
from cyclic proofs. Neil Jones \cite{2001-SCT}
uses for the same concept the words \quotationMarks{thread}
and \quotationMarks{multipath}, and states GTC, the Global Trace Condition,
as: \quotationMarks{all multipaths are in $DESC^\omega$}, where $DESC^\omega$ is the set of multipaths having some infinite descending thread.

Neil Jones terminology is probably more suitable when speaking
of programs, while trace and path appear more suitable when speaking of a proof-graph, as we do in this paper.
Definitions are the same.

We call value flow any assignment
of values in $\N$ to each $x \in X$, the assignment 
depends on a time parameter $t$. We require that the the set of instants
of time is some initial segment of the set $\N$ of natural numbers.

We repeat the same definitions more formally:

\begin{definition}[Path, Trace and Value Flow]
Assume $\setR$, $\setG$ are as above.
Assume $D = [0,h]$ for some $h \in \N$ or $D=\N$ 
is any non-empty initial segment of the order on natural numbers.
\begin{enumerate}
\item
We call any map $b:D \rightarrow \setR$ a multipath, a \emph{path} for short:
a path is any finite or infinite sequence over $\setR$ which is
\underline{composable} in
$\setG$.
\item
We call any map $\pi:D \rightarrow X$ a \emph{trace} (or a \emph{thread}): 
a trace is any finite or infinite sequence over $X$.
\item
We call any map $\theta:D,X \rightarrow \N$ a \emph{value flow}. 
We call $\theta(t,x) \in \N$ the value of $x \in X$ at time $t \in D$.
\end{enumerate}
We say that a path, trace, value flow 
is finite when $D=[0,h]$ for some $h \in \N$, that it is infinite when $D = \N$.
\end{definition}



We are interested to a condition on $\setR, \setG$
called the \quotationMarks{Global Trace Condition} (GTC), 
equivalent to the 
termination of value flows \quotationMarks{compatible}
with a path of $\setR, \setG$. GTC is expressed in term of traces.

\begin{definition}[Global Trace Condition (GTC) ]
Assume $\setR, \setG$ are as above.
Suppose $X=\{x_1, \ldots, x_n\}$, $\setR =\{(\Stat_i, \Prog_i)| i \in [1,r]\}$
and $b, \pi$ are a path of $\setR, \setG$ and a trace, all of domain $D$ 
$D=[0,h]$ for some $h \in \N$ or $D=\N$.

\begin{enumerate}

\item
We say that the trace $\pi:D \rightarrow X$ is compatible with the path
$b :D \rightarrow \setR$ on the same $D$, 
and we write $\pi \models b$, if and only if :
whenever $t, t+1 \in D$ and $b(t) = (\Stat_i,\Prog_i) \in \setR$
the two consecutive values  of the trace are related by $\setR$, namely
 $\pi(t) (\Stat_i \cup \Prog_i) \pi(t+1)$.
If furthermore we have $\pi(t) \Prog_i \pi(t+1)$,
we say that $t \in D$ is a \emph{$b$-progress point} of the trace $\pi$.

\item
$\setR,\setG$ satisfy the \emph{Global Trace Condition} (GTC) if and only if
for all infinite paths $b$ in $\setR,\setG$
there is some infinite trace $\pi:\N \rightarrow X$ such that
$\pi \models b$ and $\pi$ has infinitely many $b$-progress points.

\item
We write $\universe{(\Stat, \Prog)}
= \Stat \cup \Prog$ for the \emph{union} of the pair $(\Stat, \Prog) \in \setR$.


\item
$\setR$ satisfies \emph{injectivity} if and only if for all $i \in [1,r]$ 
the relation $\Stat_i \cup \Prog_i$ is injective.

\end{enumerate}

\end{definition}

The injectivity condition was already considered by Lee and Ben-Aram in 
\cite{2009-Lee-Aram}, Example 1.6, page 7.
Lee and Ben-Aram use the word \emph{fan-in free} for injectivity.

If formulated as above, the Global Trace condition (GTC) is equivalent to
Size Change Termination (SCT), introduced by Neil Jones \cite{2001-SCT}. 
SCT was originally used as a sufficient, and quite broad too, 
condition for termination of recursive calls. 
When we use GTC as a correctness criterion for cyclic proofs,
\emph{we add one more condition, \underline{injectivity} of the trace relations}, 
a condition which is missing in the SCT paper. 

Injectivity implies that in cyclic proofs two traces of different
variables never cross each other and significantly simplifies traces.
Instead, in the SCT condition two traces, say, $x_1 \ge x_2 > x_3$ and
$y_1 > y_2$, from two different variables $x_1$, $y_1$,
can meet into a single variable $z$. For instance, we can have:
\begin{center}
$x_1 \ge x_2 > x_3 \ge z$ 
\ \ \ and \ \ \ 
$y_1 > y_2 \ge z$
\end{center}

In this paper we study GTC in the case $\setR$ is injective, and from a
purely combinatorial viewpoint, as a termination condition. Our
result can be used for checking termination of recursive calls, with a 
faster algorithm but under less general conditions than in 
Neil Jones \cite{2001-SCT}. 
Our result could be used for translating cyclic proofs into ordinary proofs. 
Our goal is finding a termination proof using Intuitionistic Logic and First 
Order Arithmetic.

Several decidable equivalent are known for GTC. In Classical Logic, GTC for 
$\setR,\setG$ is equivalent to the fact that all value flows compatible 
with some path are finite.

We assume injectivity and we are looking for an equivalent version
of GTC which can expressed in First Order Aritmetic, and such that
the finiteness of all compatible value flows can be proved
in Intuitionistic Arithmetic.

We anticipate that when $n = \Card(X) \ge 1$, we will 
use lexicographic induction on a length $2n$ sequence 
$( m_1,r_1, \ldots, m_n,r_n)$.
For all $k \in [1,n]$, we will have $m_k \in \N \cup \{\infty\}$ and 
$r_k \in [0, {n \choose k }[$, a finite initial segment of $\N$. 



\subsection{Comparing our Terminology 
with the Terminology in Neil Jones SCT Paper}
\label{subsection-comparing}
The SCT paper uses $3$-coloring on the set $[1,n] \times [1,n]$. 
The choice of a color for a pair
$(i,j)$ represents what we know about the relation between the value of
variable $x_i$ at time $t$, and the value of variable $x_j$ at time $t+1$.
What we know about the relation can be: either $\ge$ or $>$ or 
\quotationMarks{no information at all}.

We could interpret our relations pairs
$(\Stat,\Prog) \in \setR$ as the $3$-coloring $(\Stat \setminus \Prog, \Prog, X^2 \setminus{(\Stat \cup \Prog)}$ in the sense of Neil Jones \cite{2001-SCT}. $x \Stat y$ for 
$x, y, \in X$ represents \quotationMarks{the value of $x$ is $\ge$ than the value of $y$}, provided we do not have $x \Prog y$. $x \Prog y$ represents 
\quotationMarks{the value of $x$ is $>$ than the value of $y$}. 
If we do not have $x \Stat y$ nor $x \Prog y$ we represent the fact we do
not know what is the relation between the values of $x$ and $y$.


In order to formulate SCT in our setting, we have first to recall the definition of
composition of pairs in $\setR$ and and of composable pairs in $\setG$.

%13:50 04/03/2025

\subsection{Definition of Composition of relation pairs}
\label{subsection-composition}

We define composition in such a way that if  
$(x,y), (y,z) \in X \times X$ 
are two pairs on $X$,
then $(x,z) \in X \times X$ is stationary in the composition  
$(\Stat,\Prog)(\Stat',\Prog')$ if and only if both $(x,y)$ is stationary in 
$\Stat$ and $(y,z)$ is stationary in $\Stat'$, for some $y \in X$. 
Dually, $(x,z)$ is progressing in the  composition 
$(\Stat,\Prog)(\Stat',\Prog') $ if and only if either $(x,y)$ is progressing in 
$\Prog$ or $(y,z)$ is progressing in $\Prog'$ or both are, for some $y \in X$.


\begin{definition}[States, Compositions and Transition Sequences]
Assume $X = \{1, \ldots, n\}$ and $(\Stat,\Prog)$ is a pair
of finite binary relations on $X$. 
\begin{enumerate}
\item
\emph{Composition on relation pairs}.
If $(\Stat,\Prog)$, $(\Stat',\Prog')$ are relation pairs on $X$ we define 
their composition as
$$
(\Stat,\Prog)(\Stat',\Prog') 
=
(\Stat \Stat', \Stat \Prog' \cup \Prog \Stat' \cup \Prog \Prog')
$$
\item
\emph{(Composition-Closed)}
We say that $\setR$ is \emph{composition-closed w.r.t. $\setG$},
just \emph{composition-closed} for short, if and only
whenever $(\Stat'',\Prog'')$ is the composition of
$(\Stat,\Prog), (\Stat',\Prog') \in \setR$ and 
$(\Stat,\Prog) \setG (a,b)$ and $(\Stat',\Prog') \setG (b,c)$,
then $(\Stat'',\Prog'') \in \setR$ and $(\Stat'',\Prog'') \setG (a,c)$.


\item
\emph{(Closure on composition)}
We say that $(\setR',\setG')$ is the closure by composition of $(\setR,\setG)$
if and only if:
\begin{enumerate}
\item
$\setR'$ is composition-closed w.r.t. $\setG$ and 
$\setR \subseteq \setR'$, $\setG  \subseteq \setG'$
\item
whenever
$\setR''$ is composition-closed w.r.t. $\setG''$ and 
$\setR \subseteq \setR''$, $\setG  \subseteq \setG''$, then  
$\setR' \subseteq \setR''$, $\setG'  \subseteq \setG''$.
\end{enumerate}

\end{enumerate}
\end{definition}

$\setR'$ and $\setG'$ are finite and we can compute then from $\setR$ and $\setG'$, 
because $n = \Card(X)$ is finite.
However, the cardinality of $\setR'$ could be exponential in $n$. 
From now on, \emph{we will assume that $\setR$ is composition-closed w.r.t. $\setG$}.

Remark that we defined composition of relation pairs in such a way that it is
compatible with the union operator $\universe{\cdot}$ for relation pairs. 
Indeed we have
$$
\universe{(\Stat,\Prog)(\Stat',\Prog')} 
= 
\universe{(\Stat \Stat', \Stat \Prog' \cup \Prog \Stat' \cup \Prog \Prog')}
=
\Stat \Stat' \cup \Stat \Prog' \cup \Prog \Stat' \cup \Prog \Prog'
=
(\Stat \cup \Prog) (\Stat' \cup \Prog') 
= 
\universe{(\Stat,\Prog)}\universe{(\Stat',\Prog')} 
$$
We can now define SCT.

\subsection{The Size Change Termination condition (SCT)
and the Progressing Domain Condition (PDC)}
The SCT condition says:
\quotationMarks{
\emph{all relation pairs in $\setR$ related to some $(a,a) \in G$ by $\setG$
have some finite power in which some element progresses to itself}
}. 
We write SCT as follows:

\begin{center}
\emph{
$\forall (\Stat,\Prog) \in \setR, a \in G. 
\ \ ((\Stat,\Prog) \ \setG \ (a,a)) \implies \exists n \in \N \exists x \in X. 
\ (x \mbox{ progresses to $x$ in } (\Stat,\Prog)^n)$
}
\end{center}

In previous work we proved that SCT and GTC are equivalent,
using Ramsey Theorem. 

In this paper we are looking for a more
constructive and elementary proof of an equivalent result, namely the termination
of all value flows compatible with some path of $\setR$.
The proof should be suitable to a formalization in a proof assistant.
We add the injectivity assumption because relation pairs with
injectivity are enough to describe cyclic proofs, and if we have injectivity
proving a termination result requires a weaker hypothesis PDC (see below). 
Besides, using PDC it is faster to check whether a cyclic proof is correct.

The Progressing Domain Condition (PDC) says: 
\quotationMarks{all relation pairs in $\setR$ 
have a progression from $x$ to $y$ with both elements 
included in the domain of the relation}, namely

\begin{center}
\emph{
$\forall (\Stat,\Prog) \in \setR, a \in G. 
\ \ ((\Stat,\Prog) \ \setG \ (a,a)) \implies \exists x, y \in \dom(\Stat\cup\Prog). 
\ (x \Prog y)$
}
\end{center}

PDC is weaker and it is easier to check than the SCT condition. 
To check PDC, we do not have to look for some $x$ in the progressing to itself in 
some power of $(\Stat,\Prog)$. Instead it is enough to find some $x$ progressing 
to some $y$, such that both $x$, $y$ are in the domain of $\Stat \cup \Prog$.

However, PDC has the drawback that it implies 
GTC and termination only for \emph{injective} relation pairs. 
We explain why.

\subsubsection{PDC and composition-closure do not imply GTC,
in the case of non injective family of relation pairs}
\label{subsubsection-counterexample-PDC-GTC}
Here is a counter-example $\setR_1$.
Let $X = \{x_1,x_2\}$, $\setR_1=\{(\Stat,\Prog)\}$
and $\Stat =\{(x_1,x_1)\}$, $\Prog =\{(x_2,x_1)\}$ and
$(\Stat,\Prog) \ \setG (1,1) $, 
namely, the pair $(\Stat,\Prog)$ is self-composable. 
Then $\setR_1$ is composition-closed in $\setG$ because 
$
(\Stat,\Prog)(\Stat,\Prog) = 
(\Stat\Stat,\Stat\Prog \cup \Prog\Prog \cup \Prog\Stat) = 
$(by $\Stat\Prog = \Prog\Prog = \emptyset$)$
(\Stat\Stat,\Prog\Stat) = 
(\Stat,\Prog)
$.
The unique pair $(\Stat,\Prog)$ satisfies PDC 
because $x_2 \in \dom((\Stat,\Prog))$
is progressing to  $x_1 \in \dom((\Stat,\Prog))$. However, $\setR_1$
does not satisfy GTC, because the unique infinite path in $\setR,\setG$ is 
$\pi = ((\Stat,\Prog),(\Stat,\Prog),(\Stat,\Prog),\ldots)$,
and any infinite trace compatible with $\pi$ is either
$x_1 x_1 x_1 \ldots$ or $x_2 x_1 x_1 x_1 \ldots$. Thus, any infinite
trace compatible with $\pi$ progresses \emph{at most once}. 
The implication PDC $\implies$ GTC is false because 
$\setR_1$ is not injective. Indeed, the relation 
$R = \universe{(\Stat,\Prog)} = \{(x_1,x_1), (x_2,x_1)\}$ is not injective.


Instead, we will prove that PDC+injectivity imply finiteness of value flows,
or, using Classical Logic, GTC:

\begin{theorem}[Finiteness of $\setR$-compatible transitions (weak form)]
\label{theorem-finiteness-value-flow-1}
If $\setR$ is any finite list of \emph{injective} relation pairs 
satisfying the Progressing Domain Condition (PDC) 
and composition-closed in $\setG$, 
then any value flow compatible with $\setR,\setG$ is finite.
\end{theorem}

We postpone the proof of Thm. 
\ref{theorem-finiteness-value-flow-1} above.

Before we remark that the Termination Theorem for $\setR, \setG$ 
can be reformulated as an induction schema on $\N^n$, as follows.

%%%%%%%%%%%%%%%%%%%%%%%%%%%%
% 04/03/2025 
% It is more difficult to formulate the termination theorem
% relativized to some \setG as an induction schema, we need a parameter a \in G
%%%%%%%%%%%%%%%%%%%%%%%%%%%%
%Let $\vec{x}= (x_1, \ldots, x_n)$ and $\vec{y} = (y_1, \ldots, y_n)$.
%For any $(\Stat,\Prog) \in \setR$ we define a relation $<_{(\Stat,\Prog)}$.
%We define $\vec{x}  <_{(\Stat,\Prog)} \vec{y}$ if and only if:
%$$
%\forall i,j \in [1,n]. \  
%(i \Stat j \implies x_i \le y_j) \ 
%\wedge \ 
%(i \Prog j \implies x_i < y_j)
%$$
%We define $<_\setR$ as the union 
%$\bigcup \{ <_{(\Stat,\Prog)} | (\Stat,\Prog) \in \setR\}$
%of all relations $<_{(\Stat,\Prog)}$ for $(\Stat,\Prog) \in \setR$. 
%Then Thm. \ref{theorem-finiteness-value-flow-1} can be reformulated 
%as follows. Assume PDC 
%and injectivity for $\Stat \cup \Prog$, for all $(\Stat,\Prog) \in \setR$,
%and composition-closure of $\setR$ in $\setG$.
%Then we can prove {\bf the induction schema for $<_\setR$}.
%%17:24 03/02/2025
%
%More in detail,
%using Progressing Domain Condition we will reformulate a cyclic proof
%as an ordinary proof by adding first order intuitionistic reasoning and 
%induction on $[0,\infty]$ nested $n$ times, 
%where $n = \Card(X)$ is the cardinality of the set of variables,
%together with induction on the finite sets $[1, \genfrac(){0pt}{0}{n}{k}[$ nested $n$ times,
%for $k=1, \ldots, n$.
%\\
%%%%%%%%%%%%%%%%%%%%%%%%%%%%
% END induction schema
%%%%%%%%%%%%%%%%%%%%%%%%%%%%

A last remark: in the definition of PDC we cannot simplify 
$x, y \in \dom(\Stat \cup \Prog)$ to the weaker condition
$x, y \in X$. Let us call \quotationMarks{Progressing Condition} or 
PC the resulting condition. Then PC it is weaker than PDC,
and indeed PC just says:
for all $(\Stat,\Prog) \in \setR$, all $a \in G$ if $(\Stat,\Prog) \setG (a,a)$
then $\Prog \not = \emptyset$.
If we have PC in the place of PDC then the Termination Theorem fails, even if
we assume composition-closure of $\setR$ in $\setG$ and injectivity
of $\setR$. We explain why.

\subsubsection{PC, composition-closure and injectivity does not imply GTC}
\label{subsubsection-counterexample-PC-PDC}
Here is a counter-example $\setR_2$. 
Let $X = \{x_1,x_2\}$, $\setR_2=\{(\Stat,\Prog)\}$
and $\Stat =\{(x_1,x_1)\}$, $\Prog =\{(x_1,x_2)\}$, 
namely, the pair $(\Stat,\Prog)$ is self-composable. 
Then $\setR_2$ is composition-closed because 
$
(\Stat,\Prog)(\Stat,\Prog) = 
(\Stat\Stat,\Stat\Prog \cup \Prog\Prog \cup \Prog\Stat) = 
$(by $\Prog\Prog = \Prog\Stat = \emptyset$)$
(\Stat\Stat,\Stat\Prog) = 
(\Stat,\Prog)
$.
$\setR_2$ is injective, because the relation 
$R = \universe{(\Stat,\Prog)} = \{(x_1,x_1), (x_1,x_2)\}$ is injective.
The unique coloring $(\Stat,\Prog)$ satisfies PC 
because $x_1 \in \dom((\Stat,\Prog))$
is progressing to  $x_2$. Remark that $x_2 \not \in \dom((\Stat,\Prog))$
and $x_1$ does not progress to $x_1$, therefore PDC fails while PC holds. 
$\setR_2$ does \emph{not} satisfy GTC, because the unique infinite path is 
$\pi = (R,R,R,R,\ldots)$,
and the unique infinite trace in $\pi$ is 
$x_1 x_1 x_1 \ldots$, therefore the unique infinite trace \emph{never} progresses. 
%17:24 03/02/2025

In order to prove the Termination Theorem we have to define States and
Transition chains first.




\subsection{Definition of States and Transitions chain}
\label{subsection-transition}

We define the set $\setS$ of computation states and the set $\setT$ of 
transition chains from a computation state to another one.

\begin{definition}[States, Compositions and Transition Sequences]
Assume $\setR, \setG$ are as above.
Assume $X = \{x_1, \ldots, x_n\}$ and $(\Stat,\Prog)$ is a pair
of finite binary relations on $X$. 
\begin{enumerate}

\item
We denote with $\Val$ the set of \emph{value assignment} 
$\sigma: X \rightarrow \N$ for variables. 

\item
We denote with $\setS = \setR \times \Val$ the set of \emph{computation states},
each $s=((\Stat,\Prog),\sigma)$ consists of a pair $(\Stat,\Prog)$ 
of finite binary relations on $X$ 
and a value assignment $\sigma: X \rightarrow \N$.

\item
\emph{(Transition)}
Assume $s=((\Stat,\Prog), \sigma)$, $s'=((\Stat',\Prog'), \sigma')$, $s, s' \in \setS$
are states. 
We call $(s,s')$ a \emph{transition in $\setR,\setG$}, just
a \emph{transition} for short,  if and only if
$(\Stat,\Prog)$ is composable with $(\Stat',\Prog')$ in $\setG$ and
$\forall (x,y) \in \Stat. \sigma(x) \ge \sigma'(y)$
and $\forall (x,y) \in \Prog. \sigma(x) > \sigma'(y)$.

\item
\emph{(Transition Sequence)}
A \emph{transition sequence} 
is any sequence $(s_1, \ldots, s_n) \in \setS$ of states
such that $(s_i, s_{i+1})$ is a transition in $\setR,\setG$ for all $i \in [1,n[$. 
We denote the set of transition sequences with $\setT$.
\end{enumerate}
\end{definition}


\subsection{Some Ingredient for a Proof of the Termination Theorem 
\ref{theorem-finiteness-value-flow-1}}
\label{subsection-proof-termination}
Assume that $\setR$ is a composition-closed set, that $\setR$ is
injective and that $\setR$
satisfies PDC, namely: for all $(\Stat,\Prog) \in \setR$ we have $x \Prog y$ for
some $x, y \in \dom(\Stat \cup \Prog)$.
Let $\setT$ be the set of transition sequences for $\setR, \setG$. 
We want to prove the Termination Theorem: 
all transition sequences terminate.

We will prove using Intuitionistic Logic  and First Order Arithmetic
that the tree $\setT$ of transition sequences (having root the empty sequence)
satisfies the induction schema.

To this aim, let $[0,\infty] = \N + \{\infty\}$ be $\N$ extended with an
\quotationMarks{infinite} point $\infty$, and with the relation $n < \infty$
for all $n \in \N$. We define a well-founded order $\setL$:

\begin{definition}[The well-founded order $\setL$]
\label{definition-well-founded-order-L}
Let $g = \Card(G)$. We equip the set 
$$
\setL 
\ \ \ = \ \ \ 
\Pi_{k=1}^{k=n} \ [0,\infty] \times [0,g \genfrac(){0pt}{0}{n}{k}[ 
$$
with the lexicographic ordering on the components $(m_1,r_1), \ldots, 
(m_n,r_n)$, with each pair $(m_i,r_i)$ ordered lexicographically. 
\end{definition}

%We define a map 
%$
%f: \setT \rightarrow \setL
%$ 
%and we prove
%that $f$ is increasing w.r.t. the one-step extension $>_1$ in $\setT$
%and the lexicographic ordering $<_{\lex}$ on $\setL$.
%
%Namely we will prove that if $\theta \in \setT$ and 
%$\eta \in \setT$, $\eta >_1 \theta$ then $ f(\eta) <_{\lex} f(\theta) $. 
%This implies that all transition sequences in $\setT$ terminate.

We need some more ingredients for a proof of the Termination 
Theorem \ref{theorem-finiteness-value-flow-1}. 
The first one is the following folk-lore Proposition, some sufficient condition
for an injective relation $R$ on $X$ to be a bijection.

\begin{proposition}[Injective Relations and Cardinality]
\label{proposition-injective-relation}
Assume that $R \subseteq X \times Y$ is \emph{everywhere defined} 
and is \emph{injective}. Then 
\begin{enumerate}
\item
\label{proposition-injective-relation-01}
$\Card(X) \le \Card(Y)$

\item  
\label{proposition-injective-relation-02}
if $\Card(X)=\Card(Y)$
then $R$ is a bijection between $X$ and $Y$.
\end{enumerate}
\end{proposition}

\begin{proof}
Assume $R \subseteq X \times Y$ is everywhere defined and injective. 
\begin{enumerate}
\item
%\label{proposition-injective-relation-01}
Then
for all $x \in X$ the sets $R(x) = \{ y \in Y | R(x,y) \}$ are inhabited,
because $R$ is everywhere defined, and are pairwise disjoint,
because $R$ is injective.
We deduce that $\Card(X) \le \Card(Y)$.

\item
%\label{proposition-injective-relation-02}
Now suppose that $\Card(X) = \Card(Y)$.
Then all sets $R(x)$ have cardinality $1$, 
otherwise we would have $\Card(X)<\Card(Y)$. 
Thus, $R$ is a function, and $R$ is injective by assumption. 
Now we use the well-known fact that any injection
between two \emph{finite} sets of equal cardinality is a bijection.
\end{enumerate}
\end{proof}

% 13:52 20/02/2025

Next, we also need a sufficient condition for a relation $R$ being
a bijection between $\dom(RS)$ and $\dom(S)$: $R$ should be injective
and the cardinalities of $\dom(RS)$ and $\dom(S)$ should be the same.
We also provide a sufficient condition for $R$ being
a bijection between $\dom(R)$ and $\dom(R)$:
$R$ should be injective and we should have $\dom(R)=\dom(RS)=\dom(S)$.

From now on, for any binary relation $R$ on $X$
we write $R \restr Y,Z$  
for $R \cap (Y \times Z)$, the \emph{restriction of $R$ to $Y,Z$} and 
$R \restr Y$ for $R \cap (Y \times Y)$,
the \emph{restriction of $R$ to $Y$.}

\begin{lemma}[Composition of Injective Relations and Cardinality]
\label{lemma-composition-injective-relations}
Assume that $R$,$S$ are finite binary relations on $X = \{x_1,\ldots,x_n\}$.

\begin{enumerate}
\item
\label{lemma-composition-injective-relations-01}
$\dom(R) \supseteq \dom(RS)$.

\item
\label{lemma-composition-injective-relations-02}
If $R$ is injective then 
$R \cap (\dom(RS) \times \dom(S))$ is everywhere defined and
surjective, and $\Card(\dom(RS)) \le \Card(\dom(S))$

\item
\label{lemma-composition-injective-relations-03}
If $R$ is injective and $\Card(\dom(RS)) = \Card(\dom(S))$
then $R\cap (\dom(RS) \times \dom(S))$ is a bijection.

\item
\label{lemma-composition-injective-relations-04}
If $R$ is injective and $\dom(R)=\dom(RS)=\dom(S)$ 
then $R\restr \dom(R)$ is a bijection.

\end{enumerate}
\end{lemma}

\begin{proof}
%11:52 06/02/2025

\begin{enumerate}
\item
%\label{lemma-composition-injective-relations-01}
We have to prove that $\dom(R) \supseteq \dom(RS)$. 
We assume $x \in \dom(RS)$ in order to prove that $x \in \dom(R)$.
If  $x \in \dom(RS)$ then by definition of $RS$ 
we have $(x,y) \in \dom(R)$ and $(y,z) \in \dom(S)$ for some $y,z \in X$.
We conclude that $x \in \dom(R)$, as wished.

\item
%\label{lemma-composition-injective-relations-02}
By Lemma 
\ref{proposition-injective-relation}.
\ref{proposition-injective-relation-01} 
if $R$ is injective then in order to prove that
$\Card(\dom(RS)) \le \Card(\dom(S))$ it is enough to prove two conditions
on $R \cap (\dom(RS) \times \dom(S))$:

\begin{enumerate}
\item
\emph{$R \cap (\dom(RS) \times \dom(S))$ is everywhere defined}.
For, we assume that $x \in \dom(RS)$ in order to prove that
for some $(x,y) \in R$ we have $y \in \dom(S)$.
From $x \in \dom(RS)$ we deduce that there is some $z \in X$ 
such that $(x,z) \in RS$, therefore for some $y \in X$ we have  $(x,y) \in R$ 
and $(y,z) \in S$. We conclude that $(x,y) \in R$ and  
$y \in \dom(S)$, as we wished to prove.

\item
\emph{$R \cap (\dom(RS) \times \dom(S))$ is injective}. This
is  because this relation is a subset of $R$ which is injective.
\end{enumerate}

\item
%\label{lemma-composition-injective-relations-03}
By point \ref{lemma-composition-injective-relations-02} above 
and Proposition 
\ref{proposition-injective-relation}.\ref{proposition-injective-relation-02}.

\item
%\label{lemma-composition-injective-relations-04}
Assume $\dom(R)=\dom(RS)$ and $\dom(R)=\dom(S)$ 
and $R$ injective in order to prove that $R\restr \dom(R)$ is a bijection. 
$\dom(R)=\dom(S)$ implies that $\Card(\dom(R))=\Card(\dom(S))$.
By $\Card(\dom(R))=\Card(\dom(S))$, $R$ injective
and point \ref{lemma-composition-injective-relations-03} above we deduce 
that $R\restr \dom(RS), \dom(S)$ is a bijection. The thesis now follows from
$\dom(RS) = \dom(R)$ and $\dom(S)=\dom(R)$.

\end{enumerate}
\end{proof}

%13:24 06/02/2025
%12:21 17/02/2025

Remark that the fact that $R\restr \dom(R)$ is a bijection 
does not mean that $R$ itself
is a bijection. For instance, an element $x \in \dom(R)$ can have some image
$y \not \in \dom(R)$. In this case $x$ has at least two images and $R$ is no
bijection, while $R\restr \dom(R)$ is.

The last ingredient we need is a definition $\val(Y,\sigma)$ of value for 
a subset $Y \subseteq X$ of $X$ w.r.t. a valuation 
$\sigma:X \rightarrow \N$. 
$\val(Y,\sigma)$ is the sum of all values of $\sigma$ on $X$ 
$$
	\val(Y,\sigma)  \ \ \ = \ \ \  \Sigma_{y \in Y} \sigma(y)
$$


\begin{lemma}[The value of progressing bijection is decreasing]
\label{lemma-decreasing-progressing-bijection}

Assume that $r = (\Stat,\Prog)$ is a relation pair of $X = \{x_1, \ldots, x_n\}$
and $R= \universe{r} = \Stat \cup \Prog$.
Assume $\sigma, \tau \in \Val$ are value assignments on $X$
and $Y, Z \subseteq X$. Assume that:

\begin{enumerate}
\item
$R \cap (Z \times Y)$ is a bijection 
\item
we have
\quotationMarks{weak progression of $\sigma$ to $\tau$}:
for all $i \in Y, j \in Z. ((i,j) \in R \implies \sigma(i) \ge \tau(j))$.
\item
we have \quotationMarks{strict progression of $\sigma$ to $\tau$}:
for some $i \in Y, j \in Z$ we have
$((i,j) \in R)$ and $(\sigma(i) > \tau(j))$.
\end{enumerate}

Then $\val(Y,\sigma) > \val(Z,\tau)$.

\end{lemma}


\begin{proof}
From $R \cap (Z \times Y)$ bijection we deduce that
 $\Card(Y) = \Card(Z) = p$ for some $p \in [0,n]$. 
We also deduce that $Y = \{i_1, \ldots, i_p\}$ and $Z = \{j_1, \ldots, j_p\}$
for some increasing sequences $i_1 < \ldots < i_p$
and $j_1 < \ldots < j_p$ on $[1,n]$ such that $(i_h, j_h) \in R$ for all
$h \in [1,p]$.

By weak and strict progression of $\sigma$ to $\tau$ we deduce that 
$\sigma(i_h) \ge \tau(j_h)$ for all $h \in [1,p]$
and $\sigma(i_h) > \tau(j_h)$ for some $h \in [1,p]$
We conclude that 
$
\val(Y,\sigma)                               = 
\Sigma_{h \in [1,p]} \sigma(i_h) >  
\Sigma_{h \in [1,p]} \tau(j_h)      =
\val(Z,\tau)
$,
as wished. 
\end{proof}

\section{The notion of stable-cardinality path}
We need one last ingredients before proving the Termination Theorem.
Suppose $\Card(X) = n$. Suppose $\setR=\{(\Stat_i,\Prog_i) | i \in [1,a]\}$
is any finite set of relation pairs on $X$, $\setG$ a relation between $\setR$
and $G \times G$, and $\setR$ is composition-closed in $\setG$.

We now define larger and larger set of paths, then we will prove that the
associated value flows terminate.

Assume $L:[1,l] \rightarrow \Pfin(\N \times \N)$ is a finite relation sequence,
with $R_i = L(i)$ for all $i \in [1,l]$.
Then for all $i,j \in [1,l]$, $i \le j$ we define
$R_{i,j} = \Pi_{h=i}^{h=j} R_j = R_i \ldots R_j$.

By Lemma \ref{lemma-composition-injective-relations-01} 
and $\dom(R_{i,j+1}) = \dom(R_{i,j} R_{j+1})$ we have
$\dom(R_{i,j})  \supseteq \dom(R_{i,j+1})$, that is,
the sequence of sets $\dom(R_{i,j})$
for $j \in [i,l]$ is weakly decreasing. The same holds for 
$\Card(\dom(R_{i,j}))$.
When for each $i \in [1,l]$  the sequence $\Card(\dom(R_{i,j}))$ 
is constant we say that the list $L$ has stable-cardinality domains. 
When all cardinalities are $\Card(\dom(R_{i,j}))$ are the same we
say that $L$ has constant-cardinality domains.

Here is the formal definition.

\begin{definition}
[Stable-Cardinality Domain for a sequence on $\setR$]
Assume $L:[1,l] \rightarrow \Pfin(X \times X)$
is a finite relation sequence and
$\pi:[1,l] \rightarrow \Pfin(X \times X)^2$ is
a finite sequence of relation pairs.

\begin{enumerate}

\item
$L$ has \emph{stable-cardinality domains} if and only if all $i \in [1,l]$ 
we have $\Card(\dom(R_i)) = \Card(\dom(R_{i,l}))$.

\item
$L$ has \emph{constant cardinality domains} if and only if all 
$i,j \in [1,l]$, $i \le j$ 
we have $\Card(\dom(R_{i,j})) = \Card(\dom(R_{1,l}))$.

\item 
$\pi$ has \emph{stable-cardinality domains} 
(respectively, has \emph{constant cardinality domains})
if and only if the associated relation sequence 
$\{\universe{\pi(i)} | i \in [1,l]\}$, obtained by taking the union
of each pair $\pi(i) = (\Stat_i,\Prog_i)$, has stable-cardinality domains 
(respectively, has constant cardinality domains).

\item
$\pi$ has \emph{non-progressing domains} if and only if for all 
$i, j \in [1,l]$, $i < j$ we have $\neg (x \Prog_{i,j} y)$ for all 
$x \in \dom(\pi(i))$, $y \in \dom(\pi(j))$.
\end{enumerate}

\end{definition}

We can now restate the Termination Theorem in a more detailed form, giving
an upper bound the height of the tree of all paths in various cases.
 
 \begin{theorem}[Finiteness of $\setR$-Compatible Transitions (Srong Form)]
\label{theorem-finiteness-value-flow-2}
Let $g = \Card(G)$. 
\begin{enumerate}
\item
\label{theorem-finiteness-value-flow-2-1}
The value flows on $\setR, \setG$ 
with \emph{constant-cardinality non-progressing domains} 
terminate in $< g\genfrac(){0pt}{0}{n}{k}$ steps, where $k$ is 
the common cardinality of the domains.
Therefore these value flows terminate.

\item 
\label{theorem-finiteness-value-flow-2-2}
There is a decreasing embedding $w$ from the value flows on $\setR, \setG$ 
with \emph{constant-cardinality domains} to the set 
$[0,\infty] \times [0,g\genfrac(){0pt}{0}{n}{k}[$. 
Therefore these value flows terminate.

\item
\label{theorem-finiteness-value-flow-2-3}
There is a decreasing embedding $W$ 
from the value flows on $\setR, \setG$ with 
\emph{stable-cardinality domains} to the well-order $\setL$ 
(Def. \ref{definition-well-founded-order-L}).
Therefore these value flows terminate.

\item
\label{theorem-finiteness-value-flow-2-4}
There is a decreasing embedding from \emph{all} value flows on $\setR, \setG$ 
to $L$, and equal to $WG$ for some  $G$. 
Therefore these value flows terminate.
\end{enumerate}
\end{theorem}

We postpone the proof of Thm. \ref{theorem-finiteness-value-flow-2}.
Assume $L = \{R_i | i \in [1,l\}$ and 
$R_{i,j} = \Pi_{h=i}^{h=j} R_j = R_i \ldots R_j$
for all $i,j \in [1,l]$, $i \le j$.
If $L$ has stable-cardinality domains and $i \in [1,l]$, 
then the first set $\dom(R_{i,i})$ and the last set $\dom(R_{i,l})$ in the weakly 
decreasing sequence $\{\dom(R_{i,j}) | j \in [i,l]\}$ have the same cardinality,
therefore they are equal. 
Thus, if $L$ has stable-cardinality domains then $\dom(R_{i,i}) = \dom(R_{i,j})$
for all $i,j \in [1,l]$, $i \le j$.
The next Lemma describes in details more properties of a path $\pi$ 
on a set $\setR$  with stable-cardinality domains.

\begin{lemma}
[stable-cardinality domain injective bi-coloring and the Progressing Domain condition]
\label{lemma-stable-cardinality domain-injective}
Assume that 
$\pi:[1,l] \rightarrow \Pfin(X \times X)^2$ 
is a \emph{stable-cardinality domain} finite list
of \emph{injective} pairs of relations on $X$, 
compatible with the assignment list 
$\theta:[1,l] \rightarrow (X \rightarrow \N)$.

For all $i \in [1,l]$ suppose that $R_i = \universe{\pi(i)}$. 
For all $j \in [1,l]$, $i \le j$
suppose that $R_{i,j} = \Pi_{h=i}^{h=j} R_j = R_i \ldots R_j$.

\begin{enumerate}
\item
\label{lemma-stable-cardinality domain-injective-01}
If $i,j \in [1,l]$, $i < j$ we have $\Card(\dom(R_i)) \le \Card(\dom(R_j))$.
If furthermore $\Card(\dom(R_i)) = \Card(\dom(R_j))$,
then $R_{i,j-1} \restr \dom(R_{i}), \dom(R_{j})$ is a bijection from $\dom(R_{i})$
to $\dom(R_{j})$.

\item
\label{lemma-stable-cardinality domain-injective-02}
If $i,j \in [1,l]$, $i < j$ and $\Card(\dom(R_i)) = \Card(\dom(R_j))$, then 
the value of $\theta(i)$ on any $x \in \dom(R_i)$ is $\ge$ than the value of
$\theta(j)$ on the corresponding point $y \in \dom(R_j)$ and
$\val(\dom(R_i),\theta(i)) \ge \val(\dom(R_j),\theta(j))$.

\item
\label{lemma-stable-cardinality domain-injective-03}
If $i,j \in [1,l]$, $i < j$ and $\dom(R_i)) =\dom(R_j)$, and the pair
from $i$ to $j$ is related to $(a,a)$ for some $a \in G$,
and the composition closure of $\{\pi(i) | i \in [1,l]\}$ 
satisfies the Progressing-domain condition, then 
$\val(\dom(R_i),\theta(i)) > \val(\dom(R_j),\theta(j))$.

\item
\label{lemma-stable-cardinality domain-injective-04}
Let $k \in [1,n]$
with $n = \Card(X)$.
Assume that the composition closure $\setR$ of $\setR = \{\pi(i) | i \in [1,l]\}$ 
satisfies the Progressing-domain condition PDC.
Then the number of indexes $i \in [1,l]$ such that $\Card(\dom(R_i))=k$
and $\val(\dom(R_i),\theta(i))=v$ is 
$\le g\genfrac(){0pt}{0}{n}{k}$, for all $v \in \N$.

\end{enumerate}
\end{lemma}

\begin{proof}
\begin{enumerate}

\item
%\label{lemma-stable-cardinality domain-injective-01}
If $L$ is a stable-cardinality domain sequence of injective relations then by Lemma 
\ref{lemma-composition-injective-relations}.
\ref{lemma-composition-injective-relations-02},
\ref{lemma-composition-injective-relations-03} 
we deduce that $\Card(\dom(R_i R_{i+1})) \le \Card(\dom(R_{i+1}))$
and if $\Card(\dom(R_i R_{i+1})) = \Card(\dom(R_{i+1}))$ that
$R_i \restr \dom(R_i R_{i+1}), \dom(R_{i+1})$ is a bijection.
 
For all $i \in [1,l-1]$ by stability we have 
$\dom(R_i) = \dom(R_i R_{i+1})$, therefore 
$\Card(\dom(R_i)) \le \Card(\dom(R_{i+1}))$
and if $\Card(\dom(R_i)) = \Card(\dom(R_{i+1}))$ then 
$R_i \restr \dom(R_i ), \dom(R_{i+1})$ is a bijection.

\item
%\label{lemma-stable-cardinality domain-injective-02}
By the previous point we know that 
$R_{i,j-1} \restr \dom(R_{i,j-1})$ is a bijection between $\dom(R_i)$
and $\dom(R_j)$. By compatibility between $\pi$ and $\theta$, 
the value of $\theta(i)$ on any $x \in \dom(R_i)$ is $\ge$ than the value of
$\theta(j)$ on the corresponding point $y \in \dom(R_j)$.
Thus, $\val(\dom(R_i),\theta(i)) \ge \val(\dom(R_j),\theta(j))$.

\item
%\label{lemma-stable-cardinality domain-injective-03}
Assume $\dom(R_i) = \dom(R_j)$ and the pair between $i$ and $j$
is related to $(a,a)$ for some $a \in G$. Then  $\Card(\dom(R_i)) =\Card(\dom(R_j))$
and the value of $\theta(i)$ on \emph{any} $x \in \dom(R_i)$ is 
$\ge$ than the value of
$\theta(j)$ on the corresponding point $y \in \dom(R_j)$.
By the Progressing Domain Condition, then $R_{i,j-1}$ progresses on
$\dom(R_{i,j-1})$, that is, on $\dom(R_i)$, that is, between $\dom(R_i)$ 
and $\dom(R_i) = \dom(R_j)$. 
By compatibility between $\pi$ and $\theta$, 
the value of $\theta(i)$ on \emph{some} $x \in \dom(R_i)$ is 
$>$ than the value of
$\theta(j)$ on the corresponding point $y \in \dom(R_j)$.
We conclude that $\val(\dom(R_i),\theta(i)) > \val(\dom(R_j),\theta(j))$.

%10:56 16/11/2024

\item
%\label{lemma-stable-cardinality domain-injective-04}
Assume $i,j \in [1,l]$, $i < j$ and $\Card(\dom(R_i))=\Card(\dom(R_j))=k$
and $\val(\dom(R_i),\theta(i))=\val(\dom(R_j),\theta(j))=v$.
By the previous point we deduce that $\dom(R_i) \not = \dom(R_j)$.

There are $g\genfrac(){0pt}{0}{n}{k}$ 
different cardinality $k$ subsets of an $n$ element
set $X$ with , therefore at most $\genfrac(){0pt}{0}{n}{k}$ indexes for 
$i\in [1,l]$ such that $\Card(\dom(R_i))=k$
and $\val(\dom(R_i),\theta(i))=v$ and $R_i$ is related to $(a,b)$ and $R_j$
to $(a,b')$, for some $a,b,b' \in G$.

\end{enumerate}
\end{proof}

We have now all ingredients for proving the the Termination Theorem
in strong form (\ref{theorem-finiteness-value-flow-2}).
\\

\begin{proof}
[Proof of Termination Theorem \ref{theorem-finiteness-value-flow-2}]
%13:44 03/03/2025
Assume $\setR$ is composition-closed in $\setG$. Let $g = \Card(G)$.

\begin{enumerate}

\item
%\label{theorem-finiteness-value-flow-2-1}
Lemma \ref{lemma-stable-cardinality domain-injective}.
\ref{lemma-stable-cardinality domain-injective-04} says that if 
$\setR$ is composition-closed in $\setG$,
then each constant-cardinality non-progressing domain path on $\setR$
has at most $g\genfrac(){0pt}{0}{n}{k}$ elements, that is, has
length in $[0,  g\genfrac(){0pt}{0}{n}{k}[$. This is point $1$
of Termination Theorem \ref{theorem-finiteness-value-flow-2}.

\item
%\label{theorem-finiteness-value-flow-2-2}
We assume that $\pi:[1,l] \rightarrow \setR$ is a constant-cardinality domain
path on $\setR$ and we define a decreasing map $w$ from
$(\pi,\theta) \in \setS$ to $[0,\infty] \times[0,  g\genfrac(){0pt}{0}{n}{k}[$.
We define $w(\pi,\theta) = (m,r)$ where 
\begin{enumerate}
\item
$m = \val(\pi(l),\theta(l))$ is the last value of $\val(\pi(i),\theta(i))$ 
\item
$r = g\genfrac(){0pt}{0}{n}{k} - s$, 
where $s \ge 1$ is the number of $i \in [1,l]$ 
such that $\val(\pi(i),\theta(i)) = m$.
\end{enumerate}
By point \ref{theorem-finiteness-value-flow-2-1} 
above we have $s \le g\genfrac(){0pt}{0}{n}{k}$,
therefore $r \in [0, g\genfrac(){0pt}{0}{n}{k}[$.

Now we prove that if $(\pi,\theta)$, $(\pi',\theta')$ have constant-cardinality
domains and $(\pi,\theta) <_1 (\pi',\theta')$, then $w(\pi,\theta) >
w(\pi',\theta')$ in the lexicographic ordering. Suppose $w(\pi,\theta) = (m,r)$
and $w(\pi',\theta') = (m',r')$ with $r = \genfrac(){0pt}{0}{n}{k} - s$,
$r' = g\genfrac(){0pt}{0}{n}{k} - s'$, in order to prove $(m,r) > (m',r')$.
We have $m = \val(\pi'(l),\theta') \ge m' = \val(\pi'(l+1),\theta')$.
If $m > m'$ then $(m,r) > (m',r')$. Suppose now that $m=m'$.
Then the number of of $i \in [1,l]$ such that $\val(\pi(i),\theta(i)) = m$
is greater in $\pi'$ than in $\pi$. Namely we have $s' > s'$, we deduce that
$m=m'$, $r' < r'$. Also in this case we conclude $(m,r) > (m',r')$. 

\item
%\label{theorem-finiteness-value-flow-2-2} 
We assume that $\pi:[1,l] \rightarrow \setR$ is a stable-cardinality domain
path on $\setR$, in order to define a decreasing map $W$ from
$(\pi,\theta) \in \setS$ to $\setL$ (Def. \ref{definition-well-founded-order-L}).

In each stable-cardinality domains path by Lemma 
\ref{lemma-composition-injective-relations}.
\ref{lemma-composition-injective-relations-02}
we have $\Card(\dom(R_{i,m})) = \Card(\dom(R_i R_{i+1,m})) \le
\Card(\dom(R_{i+1,m}))$. By stability we have $\dom(R_{i}) =
\dom(R_{i,m})$ and $\dom(R_{i+1}) =
\dom(R_{i+1,m})$. We deduce that $\Card(\dom(R_{i})) 
\le \Card(\dom(R_{i+1}))$: the domain cardinality in a stable-cardinality path
is weakly increasing. 

For all $k \in [1,n]$ and any domain stable $ \pi : [1,m] \rightarrow \setR $ 
we define $S(\pi,k) = \{i \in [1,m] | \Card(\dom(R_i)) = k \}$. 
We claim that each $S(\pi,k)$ is a possibly empty segment of $\pi$. 
Indeed, if $i, j \in 
S(\pi,k)$, $i \le h \le j$ then $k = \Card(\dom(R_i)) \le \Card(\dom(R_h))
\le \Card(\dom(R_j)) = k$, therefore $\Card(\dom(R_h))=k$ and $h \in S(\pi,k)$.
From $\Card(\dom(R_{i}))$ weakly increasing 
we deduce that $S(\pi,k)$ is to the left
of $S(\pi,k')$ whenever $k,k' \in [1,n]$ and $k < k'$.

Thus, we can decompose each stable-cardinality domains path 
$\pi:[1,m] \rightarrow \setR$ into
a consecutive composition $\pi_1 \conc \ldots \conc \pi_n$, 
with each $\pi_k$ the possibly empty restriction of $\pi$ to $S(\pi,k)$.
By definition, each $\pi_k$ is a \emph{constant cardinality domains} path,
indeed for each $i, j \in S(\pi,k)$, $i \le j$ we have 
$\Card(\dom(R_{i,j}))=\Card(\dom(R_i))=k$. 

Let $\theta_i$ be the restriction of $\theta$ to the segment in which is
defined $\pi_1$.
We define a list of $n$ weights $W(\pi,\theta) = (w(\pi_1, \theta_1),
\ldots, w(\pi_n, \theta_n)$ 
for a length $l$ path $\pi$ with stable-cardinality domains. 

If there is no $i \in [1,l]$ such that $\Card(\dom(\universe{\pi(i)}))=k$
then we set $w(\pi,\theta)_k = (\infty,0)$.

Now we prove that if $(\pi,\theta)$, $(\pi',\theta')$ have stable-cardinality
domains and $(\pi,\theta) <_1 (\pi',\theta')$, then $W(\pi,\theta) >
W(\pi',\theta')$ in the lexicographic ordering. 
Assume $\Card(\dom(\pi(l))) = k$. Then $\pi = \pi_1 \conc \ldots \pi_k$
and $\pi_h=\nil$ for all $h \in [k+1,n]$. We have $\pi' = \pi \conc ((\Stat,\Prog))$
for some $(\Stat,\Prog) \in \setR$ and $\Card(\dom((\Stat,\Prog))) \ge k$
by stability. of $\pi'$. Therefore $\pi'_i = \pi_i$ for all $i \in [1,k-1]$
and either $\pi'_k = \pi_k \conc ((\Stat,\Prog))$ or $\pi'_k = \pi_k$
and $\pi'_{k+1} = ((\Stat,\Prog))$.
In both cases we have $w(\pi_i,\theta_i) = w(\pi'_i,\theta'_i)_l$ for all 
$i \in [1,k-1]$. 

In the case $\pi'_k = \pi_k \conc ((\Stat,\Prog))$ 
we have  $w(\pi_k,\theta_k) > w(\pi'_k,\theta'_k)$ by point 
\ref{theorem-finiteness-value-flow-2-2}  above. 
We conclude $W(\pi,\theta) > W(\pi',\theta')$.

In the case $\pi'_k = \pi_k$ and $\pi'_{k+1} = ((\Stat,\Prog))$
we have $w(\pi_k,\theta_k) = w(\pi'_k,\theta'_k)$ and 
$w(\pi_{k+1},\theta_{k+1}) = (\infty,0) > w(\pi'_{k+1},\theta'_{k+1})$.
We conclude $W(\pi,\theta) > W(\pi',\theta')$.

\item
We assume that $\pi:[1,l] \rightarrow \setR$ is \emph{any}
path on $\setR$ and we define a map $G$ from 
$(\pi,\theta) \in \setS$ to $G(\pi,\theta) = (\pi',\theta') \in \setS$, 
with $\pi'$ paths with stable-cardinality domains, in such a way that
$WG$ is decreasing. We will conclude that
\emph{all} value flows terminate. 

%16:29 03/03/2025

Assume $L=(c_1, \ldots, c_n)$ is any path on $\setR$ with $n\ge 1$. 
We define a re-grouping $M$ of $L$ as any path
$
M = 
(c_1\ldots c_{a_1},  
c_{a_1+1}\ldots c_{a_2},
\ldots,
 c_{a_{m-1}+1}\ldots c_{n})
$ 
of composition of 2-colorings taken from $L$, 
 with the same ordering as in $L$,
 and whose product is the same as $L$: $\Pi M = \Pi L$.
 $M$ is obtained by partitioning $L$ into $m \ge 1$ adjacent segments 
$[a_1,a_2-1], [a_2,a_3-1], \ldots, [a_{m-1},a_m - 1],[a_m,n]$,
with $a_1=1$, 
then composing all $c_{a_{i}},\ldots, c_{a_{i+1}-1}$ in the same segment
 into a single coloring $d_i = c_{a_{i}} \ldots c_{a_{i+1}-1}$.
 Say: we can group $L=(c_1,c_2,c_3,c_4,c_5,c_6)$ into
 $M=(c_1c_2,c_3,c_4,c_5,c_6)$, with $d_1 = c_1c_2$,
 $d_2 = c_3c_4$, $d_3=c_5$, $d_4=_6$ and in many more ways.
We call $a=a_1, \ldots, a_m$
the adjacence sequence. If $n \ge 1$ there are $2^{n-1}$ coloring,
exactly the number of subsets $\{a_1,\ldots, a_m\}$ of $[1,n]$ including 
 $a_1 = 1$.

Formally, we define re-grouping and the adjacence sequence 
by induction on $n$. 
\begin{enumerate}
\item
If $n=0$ then $M=()$
is the only re-grouping of $L=()$ with adjacence sequence $a=()$.
\item
Suppose $n \ge 1$ and we defined re-grouping and adjacence for all lists
of length $\le n-1$ and $L=L_1 \star L_2 \star (c_n)$ 
is any decomposition of $L$ into two possibly empty 
consecutive sublists $L_1$, $L_2$, plus one element. 
Assume $M$ is any re-grouping of $L_1$ with adjacence sequence $a$.
Assume $d$ is the composition of all pairs in $L_2 \star (c_n)$.
Then $M \star (d)$ is one of the re-groupings of $L$, the re-grouping 
with adjacence sequence $a \star (\length(L_1)+1)$.
\end{enumerate}
%15:18 21/02/2025

Assume $(\pi,\theta)$ 
is any compatible pair bi-coloring sequence/assignment sequence
with $\pi:[1,n] \rightarrow \setR$ and $\theta:[1,n],X\rightarrow \N$.
We define a map $G(\pi,\theta) = (q,\eta)$
returning a compatible pair bi-coloring sequence/assignment sequence,
such that for some increasing map $a:[1,m] \rightarrow [1,n]$ we have:

\begin{enumerate}
\item 
$q:[1,m] \rightarrow \setR$ 
is a re-grouping of $\pi$ with adjacence sequence $a$,
\item
$\eta=\theta \comp a$, namely $\eta(t,x) = \theta(a(t),x)$ for all
 $t \in [1,m]$, $x \in X$. 
\item
$q$ has stable-cardinality domains and  
\item
$(\pi,\theta) <_1 (\pi',\theta')$ implies $W(G(\pi,\theta)) >_\lex W(G(\pi',\theta'))$
in the lexicographic ordering.
\end{enumerate}

%17:00 03/03/2025

Assume $(\pi,\theta) <_1 (\pi',\theta')$ and  $G(\pi,\theta) = (q,\eta)$, 
$G(\pi',\theta') = (q',\eta')$. Suppose $q= q_1 \conc \ldots \conc q_n$,
$q' = q'_1 \conc \ldots \conc q'_n$ are decomposition with each $q_k, q'_k$
 a path  with constant-cardinality domains, 
 whose domains all have cardinality $k$.

%16:42 21/02/2025
Suppose $\pi' = \pi \star (\Stat,\Prog)$ and $\theta' = \theta \star \sigma$
and $G(\pi,\theta) = (q,\eta)$ with $\length(q) = \length(\eta) = m$ in order to
define $G(\pi',\theta') = (q',\eta')$, sequences of length $m'$. 
We distinguish two cases.

\begin{enumerate}
\item
Assume that $q \star (\Stat,\Prog)$ is stable. 
Then we set $q' = q \star ((\Stat,\Prog))$ and $\eta' = \eta \star (n+1)$,
two sequences of length $l' = l+1$. 

\item
Assume that $q \star (\Stat,\Prog)$ is \emph{not} stable. 
Then we take the first $l' \in [1,l]$ such that 
$
\dom(q(l') \ldots q(l))
\supset 
\dom(q(l') \ldots q(l) (\Stat,\Prog))
$.
We define $q'(i) =  q(i)$ for all $i \in [1,l'-1]$ and
$q'(l') = q(l') \ldots q(l)  (\Stat,\Prog)$.
We define $\eta'(i) = \eta(i)$ for all $i \in [1,l']$. 
\end{enumerate}
This is $(\pi',\theta')=G(\pi,\eta)$, a path with stable-cardinality domains. 

In the first case of definition of $G$ we have 
$q' = q \star ((\Stat,\Prog))$ and $\eta' = \eta \star (n+1)$, namely
$(q',\eta') >_1 (q,\eta)$, therefore $W(G(\pi,\theta)) >_\lex W(G(\pi',\theta'))$
by point \ref{theorem-finiteness-value-flow-2-3} 
above and $(\pi',\theta')$ stable.

In the second case of definition of $G$ we have: 
for $k = \Card(\dom(q'(l')))$ we have $q_1 = q'_1$, \ldots, $q_{k-1} = q'_{k-1}$
and $q_k <_1 q'_k$ and $q' = q'_1 \conc \ldots \conc q'_k$. 
By point \ref{theorem-finiteness-value-flow-2-2} above, 
this implies that we have $W(G(\pi,\theta)) >_\lex W(G(\pi',\theta'))$.

\end{enumerate}
\end{proof}

This ends the proof of Termination Theorem 
(\ref{theorem-finiteness-value-flow-1}).

\begin{thebibliography}{9}

\bibitem{2001-SCT}
Chin Soon Lee and Neil D. Jones and Amir M. Ben-Amram,
2001, January,
The Size-Change Principle for Program Termination,
in ACM SIGPLAN Notices, Vol. 36 (3), 
doi 10.1145/360204.360210

\bibitem{2009-Lee-Aram}
Amir M. Ben-Amram and Chin Soon Lee,
Ranking Functions for Size-Change Termination II,
in Logical Methods in Computer Science
Vol. 5 (2:8) 2009, pp. 1–29

\end{thebibliography}


\newpage

\section{Appendix: PDC and SCT are Equivalent Given Injectivity}
The goal of this section is to compare PDC with the GTC with and 
without injectivity. Recal that we only consider $\setR, \setG$ such that 
\emph{$\setR$ is composition-closed in $\setG$}.
We Claim:
\begin{enumerate}
\item
\label{property-1}
We prove that SCT, GTC are equivalent with PDC given injectivity:
$$
\forall \setR, \setG \mbox{ comp.-closed}.  \ \ 
\setR \mbox{ injective} 
\implies 
(SCT \Leftrightarrow PDC) \wedge 
(GTC \Leftrightarrow PDC)
$$
\item
\label{property-2}
Instead, in general we only have left-to-rigth implication
$$
\forall \setR, \setG \mbox{ comp.-closed}.  \ \ 
(SCT \Rightarrow PDC) \wedge 
(GTC \Rightarrow PDC)
$$
\item
\label{property-3}
while the opposite implication can fail:
$$
\exists \setR, \setG  \mbox{ comp.-closed}. \ \ 
\neg(\setR \mbox{ injective}) \wedge 
\neg SCT \wedge 
\neg GTC \wedge 
PDC
$$
\end{enumerate}

The implication GTC $\Rightarrow$ PDC follows from the fact that
for all self-composable $(\Stat,\Prog) \in \setR$ 
there is some infinitely progressing trace $z_1,z_2,z_3,\ldots$
in the path $R,R,R,R,...$ where $R = \universe{(\Stat,\Prog)}$. 
In any progressing point we have $z_i \Prog z_{i+1}$
and $z_{i+1} R z_{i+2}$. We conclude that $z_i \Prog z_{i+1}$
and $z_i, z_{i+1} \in \dom(R)$ for some $z_i, z_{i+1}$, that is,
PDC.
\\

In order to prove SCT $\Rightarrow$ PDC, we recall 
that Neil Jones in \cite{SCT-2001} 
proved the equivalence of SCT and GTC with only
assumption the composition-closure: 
$$
\forall \setR, \setG \mbox{ comp.-closed}.  \ \ 
(SCT \Leftrightarrow GTC)
$$
Thus, from GTC $\Rightarrow$ PDC just proved we deduce
SCT $\Rightarrow$ PDC. We proved the property \ref{property-2}.
\\

Now we prove the property \ref{property-1}.
We first prove that if $\setR$ is injective then
PCD $\Rightarrow$ SCT. If will follow
if $\setR$ is injective then PDC $\Rightarrow$ GTC. We just
proved the opposite implications.

So it is enough to prove: ($\setR$ is injective, PCD $\Rightarrow$ SCT).
By Ramsey Theorem, for all binary $\setG$-self-composable 
relations $R$ on $X = \{x_1, \ldots, x_n\}$
there is $m \in \N$, $m \ge 1$ such that $R^m$ is idempotent, namely
such that $R^m = R^m R^m$. If we $R$ is injective, then $R^m$ is 
is the identity on $\dom(R^m)$, as the next Lemma says.


\begin{proposition}[Idempotent Injective Relations
and the Global Trace Condition]
\label{proposition-idempotent-relations}
Assume a set $\setR$ of \emph{injective} relation pairs. 
Assume $(\Stat,\Prog) \in \setR$ is self-composable, 
$R = \universe{(\Stat,\Prog)}$ and $RR=R$. 
Then:
\begin{enumerate}

\item
$R \restr \dom(R) $ is a bijection

\item
$R \restr \dom(R) $ is the identity

\item
If PDC holds, then some $x \in \dom(R)$ is progressing to itself and all
$x \in \dom(R)$ are weakly progressing to themselves.
\end{enumerate}
\end{proposition}

\begin{proof}
\begin{enumerate}
\item
We have to prove that \emph{$R \restr \dom(R) $ is a bijection}. 
Let $\dom(R) = \{y_1, \ldots, y_m\}$. Assume
$i \in [1,m]$, $y_i \in \dom(R)$: 
from $R=RR$ we deduce that $y_i \in \dom(RR)$, therefore
there are $(y_i,y) \in R$ and $(y,z) \in R$. From $(y,z) \in R$ 
we obtain $y \in \dom(R)$. 
Let $C_i = \{y \in \dom(R) |(y_i,y) \in R\}$ and $c_i = \Card(C_i)$. 
By $(y_i,y) \in R$ we have $c_i \ge 1$, hence $\Sigma_{i \in [1,m]} c_i \ge m$. 
By injectivity of $R$ if  $C_i \cap C_j \not = \emptyset$ then $i=j$, therefore 
$\{C_i | i \in [1,m]\}$ are a partition of some subset of $\dom(R)$
and we have $\Sigma_{i \in [1,m]} c_i \le m$. 
We deduce that $\Sigma_{i \in [1,m]} c_i  = m$,
and from it that $c_i=1$ for all $i \in [1,m]$. Thus, 
$R \restr \dom(R) $ is an injection on a finite set $\dom(R)$.
By the characterization of bijections on finite sets we 
conclude that $R \restr \dom(R) $ is a bijection.

\item
In the point above we proved that $R \restr \dom(R) $ is 
a bijection and we have to prove that \emph{it is an identity}. Any bijection
of the finite set $\dom(R)$ is an union of pairwise disjoint cycles of the form 
$D = \{z_1, \ldots, z_p\}$ for some $p \ge 1$. We have $z_i R z_j$
if and only if $j = i+1 \mod p$ and we have
$z_i RR z_j$ if and only if $j = i+2 \mod p$.
From $R=RR$ and $z_1 R z_2$ we deduce $z_1 RR z_2$, then by choosing
$i=1$, $j=2$ that $2=1+2 \mod p$, and eventually that $0=1 \mod p$. 
We conclude that $p=1$. Thus, all cycles in
$R \restr \dom(R) $ are unary and therefore
$R \restr \dom(R) $ is the identity.

\item

\begin{enumerate}

\item 
We prove that \emph{for some $x \in \dom(R)$, $x$ is progressing to itself}.
By PDC there are some $y, z \in \dom(R)$ such that $y$ is progressing to $z$.
From $R \restr \dom(R) $ identity (the previous point)
we deduce that $y=z$ for some $x \in \dom(R)$. 
Thus, some $y \in \dom(R)$ is progressing to itself.

\item 
We prove that \emph{for all $x \in \dom(R)$, $x$ is weakly progressing to itself}.
From $R \restr \dom(R)$ identity (the previous point)
we deduce that $x R x$, either $x$ is progressing to $x$ or 
$x$ is weakly progressing to $x$. In both cases $x$ is 
is weakly progressing to $x$.
\end{enumerate}
\end{enumerate}
\end{proof}

We already said that for all finite relations $R$ 
there is $m \in \N$, $m \ge 1$ such that $R^m = R^m R^m$. By the previous
Lemma, if $\setR$ is injective, 
$(\Stat,\Prog) \in \setR$ is self-composable  in $\setG$, 
we deduce that some $x$ in progressing to itself 
in $(\Stat,\Prog)^m$. This proves SCT, hence GTC.
\\

The last property, namely
$
\exists \setR, \setG  \mbox{ comp.-closed}. \ \ 
\neg(\setR \mbox{ injective}) \wedge 
\neg SCT \wedge 
\neg GTC \wedge 
PDC
$,
was already proved with the counter-example $\setR_2$
in \S \ \ref{subsubsection-counterexample-PDC-GTC}.


\section{Appendix: Comparing Global Trace Condition (GTC) 
with Termination of Computations (TC)}
In this section we prove that the Global Trace Condition is equivalent to
the Termination of all possible computation flows, from now on, just TC
for short.

Using Classical Logic, 
it is immediate that GTC $\implies$ TC. Indeed, in any infinite computation 
flow there is some variable $x \in X$ with some infinitely progressig trace. 
This implies that the value of $x$ is infinitely decreasing in the domain $\N$,
and this cannot be. 


The converse TC $\implies$ GTC is true but not immediate. Using Classical 
Logic, we prove it in the form $\neg$ GTC $\implies$ $\neg$ TC. Actually
we prove a stronger result: if GTC is false, there is some \emph{constant}
infinite transition $s, s, s, \ldots$ for some $s = (\Stat,\Prog),\sigma) \in \setS$,
in particular Termination of Computations (TC) fails.

\begin{proof}[$\neg$ GTC $\implies$ $\neg$ TC]
Assume $\neg$ GTC. By Neil Jones Thm we deduce $\neg$ SCT, namely
that there is some self-composable 
$(\Stat,\Prog) \in \setR$ such that for no $m \in \N$, $m>0$, 
$x \in X$ we have $x$ progressing to itself in $(\Stat,\Prog)^m$.

We will define some \emph{constant}infinite transition 
$\theta = (s, s, s, \ldots)$ for some 
state $s = ((\Stat,\Prog),\sigma) \in \setS$.

We say that a trace $\tau:[1,t] \rightarrow X$ 
is irredundant if and only if all its elements are 
pairwise distinct, namely: $\forall i,j \in [1,t]. \tau(i) = \tau(j) \implies i=j$. 
An irredundant trace has at most $n$ elements, where  $n = \Card(X)$, and
progresses at most $n-1$ times, because the last element of the trace
does not progress, having no next element.
\\

Let $\pi = (\Stat,\Prog), (\Stat,\Prog), \ldots$
be the infinite constant path repeating the self-composable pair $(\Stat,\Prog)$.
We set $\sigma(x)=$ the maximum number of progressions in
an irredundant trace $\tau$, starting
 from $x$ and compatible with $\pi$. The maximum
exists because an irredundant trace has at most $n$ elements,
therefore there are finitely many irredundant traces from $x$. In fact,
we also know that $\sigma(x) \le n-1$.
\\

Now we prove that $\theta = (s, s, s, \ldots)$ is a constant infinite transition.
We have to prove that for all $x,y \in X$ if $x \Stat y$ then $\sigma(x) \ge
\sigma(y)$ and if $x \Prog y$ then $\sigma(x) >
\sigma(y)$, namely $\sigma(x) \ge \sigma(y)+1$. 
Assume $\sigma(x) = a$ and $\sigma(y) = b$.  
By definition of $\sigma$ there is some irredundant trace $\tau$ from $y$ 
whose progressing points in $\pi$ are in number of $b$. 
We have to prove $a \ge b$ if $x \Stat y$ and $a \ge b+1$ if $x \Prog y$.
\\

For, we define a trace $\rho = (x) \conc \tau$.
$\rho$ could be not irredundant, however we can repeatedly replace each
cycle included in $\rho$ with its first element. 
Each time we obtain a shorter trace from $x$, eventually we obtain an irredundant
trace $\rho'$. All the cycles we removed are not progressing,
otherwise there would be some power of $(\Stat,\Prog)$ and some $z \in X$
 progressing to itself. Therefore whenever we remove a cycle from$\rho$
 we preserve the number of progressions. This number is $b$ if $x \Stat y$
 and is $b+1$ if $x \Prog y$. By definition of $\sigma(x) = a$ we deduce
 $a \ge b$ if $x \Stat y$ and $a \ge b+1$ if $x \Prog y$, as wished.
\end{proof}

Remark that in fact we proved a stronger property: either we have SCT and
GTC and TC, or we have $\neg$ SCT and $\neg$ GTC and $\neg$ TC. In the
second case we also have an infinite constant transition $\theta = (s, s, s, \ldots)$
in which all traces progresses at most $n-1$ times, with $n=\Card(X)$.

\end{document}

